% -------------------------------------------------------------------------------------
\paragraph{本章小结}
沿 $z$ 轴的推动把光锥分量乘以 $\ee^{\pm\eta}$、把好/坏旋量分量乘以 $\ee^{\pm\eta/2}$（\cref{lem:prelim-boost}），这是幂次计数的出发点。Wilson 线是规范协变的平行移动算符（\cref{prop:prelim-wilson-covariance}），其动量空间规则由一条统一规则给出，未来与过去指向的线相差 $2\pi\ii\,\delta(n\cdot k)$（\cref{prop:prelim-wilson-rules}）；程函恒等式在绝热正规化下精确成立（\cref{lem:prelim-eikonal-identity}），并把高能部分子与软胶子的耦合等同于 Wilson 线（\cref{ex:prelim-eikonal-two}）。维数正规化的每个极点都要标明来源，无标度积分为零（\cref{lem:prelim-scaleless,ex:prelim-scaleless}）。加号分布与 $(1-z)^{-1-a\varepsilon}$ 的展开组织了单圈实辐射（\cref{lem:prelim-plus-expansion}），Mellin 变换把卷积化为乘积（\cref{lem:prelim-mellin}），矩只把 $x$ 空间的对象确定为分布（\cref{rem:prelim-moment-tools}）。切割规则、物理幺正性与光学定理（\cref{thm:prelim-cutting-rules,prop:prelim-unitarity,prop:prelim-optical}）把截面与前向振幅的虚部联系起来；无质量维数正规化中的割图规则是全阶论证（\cref{rem:prelim-cutting-dr}，\tagA）。

\paragraph{这些工具在后文的用途}
光锥坐标与推动引理用于第~4~章的区域标度与幂次计数（\cref{def:ir-scalings}、\cref{prop:ir-power-counting}）。Wilson 线出现在部分子分布的定义（\cref{def:pdf-quark}）、尖点反常维度（\cref{def:sudakov-cusp}）、共线与软近似（\cref{def:method-collinear-approx}、\cref{def:method-soft-approx}）、多胶子 Ward 恒等式（\cref{lem:method-ward-multi}）与 TMD 的定义（\cref{def:tmd-unsub}）中。维数正规化的记账用于\cref{prop:sudakov-oneloop}与\cref{prop:pdf-partonic-dr}，$\gamma_5$ 的约定用于\cref{thm:dis-bjorken}与\cref{thm:coll-g1-moments}。加号分布与 Mellin 变换用于分裂函数（\cref{prop:pdf-lo-kernels}）、次领头阶系数函数（\cref{prop:disall-nlo}、\cref{prop:dy-nlo}）与阈值区（\cref{prop:thr-oneloop}）；矩问题的工具用于\cref{cor:dis-weak-xspace}、\cref{rem:dis-nonuniform}与\cref{rem:pdf-ns-plusminus}。割图与光学定理是\cref{thm:ee-R-finite}、\cref{prop:dis-dispersion}、\cref{prop:method-fsi}与\cref{lem:method-soft-unitarity}的基础，其光锥形式在附录~O 中证明。

\begin{table}[htbp]
\centering
\caption{本章记号速查}\label{tab:prelim-symbols}
\footnotesize
\renewcommand{\arraystretch}{1.05}
\begin{tabularx}{\textwidth}{@{}l X l@{}}
\toprule
记号 & 含义 & 定义处 \\
\midrule
$v^\pm$，$\mathbf v_{\mathrm T}$，$y$ & 光锥分量 $(v^0\pm v^3)/\sqrt2$，横向分量，快度 $\frac12\ln(k^+/k^-)$ & \cref{def:prelim-lightfront} \\
$n_A$，$n_B$ & 类光矢量，$n_A\cdot v=v^-$，$n_B\cdot v=v^+$，$n_A\cdot n_B=1$ & \cref{def:prelim-lightfront} \\
$P_\pm$ & 好/坏分量的投影 $P_+=\frac12\gamma^-\gamma^+$，$P_-=\frac12\gamma^+\gamma^-$ & \cref{lem:prelim-boost} \\
$\lambda$ & 幂次计数参数 $m/Q$；$A$-共线标度 $Q(1,\lambda^2,\lambda)$ & \cref{ex:prelim-boost-scaling} \\
$W[b,a]$ & Wilson 线 $P\exp(\ii g\int_a^b\dd z\cdot A)$，较晚的点在左 & \cref{def:prelim-wilson} \\
$W[\infty n,0]$，$W[0,-\infty n]$ & 未来指向、过去指向的半无限线 & \cref{def:prelim-wilson} \\
$W^{(A)}_{ab}$，$\tilde A(k)$ & 伴随 Wilson 线 $2\Tr[t^aWt^bW^\dagger]$；胶子场的 Fourier 分量（$k$ 流入） & \cref{def:prelim-wilson} \\
$\epsilon$ & 绝热参数（与维数正规化的 $\varepsilon$ 无关） & \cref{def:prelim-wilson} \\
$\varepsilon_{\mathrm{UV}}$，$\varepsilon_{\mathrm{IR}}$ & 维数正规化极点的来源标注 & \cref{def:prelim-dr-conventions} \\
$\gamma_5$，$\epsilon^{\mu\nu\rho\sigma}$ & $\gamma_5=\ii\gamma^0\gamma^1\gamma^2\gamma^3$，$\epsilon^{0123}=+1$ & \cref{def:prelim-gamma5} \\
$g_\perp^{\mu\nu}$，$\epsilon_\perp^{\mu\nu}$ & 横向度规与横向反对称张量，$\epsilon_\perp^{12}=+1$ & \cref{def:prelim-gamma5} \\
$[g]_+$，$\mathcal D_k$ & 加号分布，$\mathcal D_k=[\ln^k(1-z)/(1-z)]_+$ & \cref{def:prelim-plus} \\
$\tilde f(N)$，$\otimes$ & Mellin 变换 $\int_0^1z^{N-1}f$，Mellin 卷积 & \cref{def:prelim-mellin} \\
$S_k(N)$，$\bar N$ & 调和和，$\bar N=N\ee^{\gamma_E}$ & \cref{def:prelim-mellin} \\
$\Pi(q^2)$，$R(s)$ & 流关联函数，$R=12\pi\Imag\Pi(s+\ii0)$ & \cref{prop:prelim-optical} \\
$W^{\mu\nu}$，$T^{\mu\nu}$ & 强子张量与前向 Compton 振幅，$W^{\mu\nu}=\frac1{2\pi}\Imag T^{\mu\nu}$ & \cref{prop:prelim-optical} \\
$\Imag X^{\mu\nu}$ & 吸收部分 $\frac1{2\ii}[X^{\mu\nu}-(X^{\nu\mu})^*]$ & \cref{prop:prelim-optical} \\
\bottomrule
\end{tabularx}
\end{table}
